% =====================================================================
\section{Methodology}
\label{sec:method}
% =====================================================================
% Every statement in this section was checked against the repository
% (file paths given in comments). Anything that could not be verified is
% marked \TODO. Do not add details from memory.

% ---------------------------------------------------------------------
\subsection{Overview and Scope}
\label{sec:method:overview}
% ---------------------------------------------------------------------
The pipeline has five stages: (i) a fixed YOLO detector is run on UAV
orthomosaic patches; (ii) detections with confidence $\geq 0.5$ become
verification samples, one per box; (iii) each sample is rendered under five
input conditions (\Acond{1}--\Acond{5}) with pre-generated prompts; (iv) each
VLM returns a JSON verdict that is parsed by a shared parser; (v) verdicts are
analyzed (a) without any reference labels (behavioral analysis), (b) against
LabelMe annotations through IoU matching (annotation alignment), and (c) against
a human semantic review of a subset (semantic audit). LabelMe annotations are
used only for evaluation and are never shown to the VLM.
% Source: README.md "Overview"; docs/ARCHITECTURE.md; docs/FRAMEWORK_FREEZE.md

\begin{figure}[t]
  \centering
  % TODO: AUTO-GENERATE FROM RESULTS
  % Candidate source: README mermaid flowchart (redraw as vector figure).
  \FIGPLACEHOLDER{Pipeline diagram: UAV patches $\rightarrow$ frozen YOLO detections $\rightarrow$ one sample per box $\rightarrow$ A1--A5 builder $\rightarrow$ VLM adapter $\rightarrow$ shared parser $\rightarrow$ \{R, U, Ur\} $\rightarrow$ behavioral / alignment / semantic analyses. LabelMe enters only at evaluation.}
  \caption{Verification pipeline. Blue components are frozen and shared by all models; only the VLM adapter is model-specific.}
  \label{fig:pipeline}
\end{figure}

% ---------------------------------------------------------------------
\subsection{Frozen Protocol versus Later Analyses}
\label{sec:method:freeze}
% ---------------------------------------------------------------------
The verification framework was frozen in July 2026
(\texttt{docs/FRAMEWORK\_FREEZE.md}). The frozen components are fixed across all
models: verification dataset and sample identifiers, \Acond{1}--\Acond{5}
definitions and image files, prompt text, output label set, JSON output schema,
response parser, IoU threshold (0.5), greedy one-to-one matching, exclusion of
\U{} from binary metrics, metric formulas, resume convention, and treatment of
failed outputs (recorded as \texttt{parse\_error} with an empty decision; never
coerced). What may vary across models is the checkpoint, vision encoder,
processor/tokenizer, chat-template wrapping of the shared prompt text,
generation settings (e.g., dtype), and model-specific raw-response cleanup.

Several elements were added \emph{after} the freeze. They change how stored
predictions are scored or interpreted, but no stored prediction was modified and
no inference was rerun for any of them (\cref{tab:freeze_vs_later}).

\begin{table}[t]
  \centering
  \small
  \caption{Frozen experimental protocol versus later amendments and analyses. Dates are from repository documentation.}
  \label{tab:freeze_vs_later}
  \begin{tabularx}{\linewidth}{@{}>{\raggedright\arraybackslash}p{0.34\linewidth} >{\raggedright\arraybackslash}p{0.17\linewidth} >{\raggedright\arraybackslash}X@{}}
    \toprule
    \textbf{Element} & \textbf{Status} & \textbf{Effect / source} \\
    \midrule
    Detections, sample IDs, \Acond{1}--\Acond{5} images and prompts, parser, schema, decoding & Frozen (Jul 2026) & Shared by all models. \texttt{FRAMEWORK\_FREEZE.md} \\
    IoU $\geq 0.5$, greedy one-to-one matching, \U{} excluded & Frozen (Jul 2026) & \texttt{EVALUATION\_PROTOCOL.md} \\
    Protocol v2 palm-label normalization & Amendment (2026-09-27) & Case-insensitive \texttt{palm} label; re-scored all stored predictions; changed reference labels only (\cref{sec:method:labelme}) \\
    Construct reframing of Protocol v2 & Interpretation (2026-10-01) & Matched/unmatched read as \emph{annotation alignment}, not semantic truth; no number changed \\
    Human semantic audit & Later analysis (2026-10-01) & 638 \unmatched{} detections + 400-item lower-confidence pilot (\cref{sec:method:semantic}) \\
    Ministral~3 raw-response normalizer & Amendment (2026-10-01) & Escapes control characters inside JSON strings, Ministral~3 keys only; no stored prediction re-parsed \\
    GT-independent \Acond{1}--\Acond{5} behavior analysis & Later analysis & \texttt{scripts/analysis/a1\_a5\_behavior\_analysis.py} \\
    Detection-property behavior analysis & Later analysis & \texttt{scripts/analysis/detection\_property\_behavior\_analysis.py} \\
    IoU-threshold sensitivity & Later analysis & \texttt{scripts/analysis/iou\_sensitivity\_analysis.py}; canonical threshold unchanged \\
    Model-scaling campaign and paired-probe qualification & New runs under the frozen protocol (Oct 2026) & \texttt{docs/MODEL\_EXPANSION\_CAMPAIGN.md} \\
    \bottomrule
  \end{tabularx}
\end{table}

% ---------------------------------------------------------------------
\subsection{Imagery and Candidate Detection Construction}
\label{sec:method:detections}
% ---------------------------------------------------------------------
\paragraph{Imagery.} Inputs are UAV orthomosaic patches stored as PNG files
together with LabelMe JSON annotations (880 annotation files).
% Source: src/paths.py RAW_PATCHES_ROOT; docs/EXPERIMENT_RESULTS_CANONICAL.md §3
\TODO{Study site, acquisition platform/sensor, flight altitude, ground sampling distance, orthomosaic generation, and patch-tiling procedure. None of these are documented in the repository.}
\TODO{Patch dimensions: the one patch checked (source of \texttt{sample\_000001}) is $912\times912$ px; verify for all patches.}

\paragraph{Detector.} Candidates come from a single YOLO11x model
(Ultralytics) with weights \texttt{yolo11x\_palm\_new/weights/best.pt}. The
detector is frozen and is not trained or modified in this work.
% Source: src/paths.py YOLO_MODEL_PATH; scripts/run_full_inference.py
\TODO{Detector training data, training procedure, and whether the training set overlaps the 880 evaluation patches. Weights come from a separate training directory not tracked in this repository.}

Inference was run on all raw patches with confidence threshold $0.001$ and at
most 300 detections per image, producing 123{,}405 detections in total.
% Source: scripts/run_full_inference.py (CONF_THRESHOLD=0.001, MAX_DET=300);
%         docs/EXPERIMENT_RESULTS_CANONICAL.md §3 (123,405)
\TODO{Inference image size and NMS settings (Ultralytics defaults were used; record the exact values from the Ultralytics version used).}

\paragraph{Verification cohort.} Detections with YOLO confidence
$\geq 0.5$ were retained, giving $N = 5{,}747$ detections on 870 images
(confidence range 0.500--0.924, median 0.834). Each retained detection becomes
one verification sample. Samples are sorted by image identifier and then by
descending confidence, and assigned sequential identifiers
\texttt{sample\_000001}--\texttt{sample\_005747}.
% Source: src/preprocessing/verification_dataset.py (iter_verification_samples,
%         group_samples_by_image); scripts/pipeline/generate_verification_dataset.py
%         (--confidence default 0.5); outputs/verification_dataset/index.csv
Consequently the first $k$ samples (used for @100/@1000 qualification runs) are
\emph{not} a random sample of the cohort.
\TODO{Justify the 0.5 confidence threshold (not documented beyond being the default).}

For each sample the stored metadata are the box $(x, y, w, h)$, confidence,
width, height, area, and center coordinates. Box widths range from 26 to 357~px
(median 150~px).
% Source: build_detection_metadata(); index.csv describe()

% ---------------------------------------------------------------------
\subsection{LabelMe Reference Annotations}
\label{sec:method:labelme}
% ---------------------------------------------------------------------
Reference annotations are LabelMe shapes. Under Evaluation Protocol~v2 a shape
is a palm iff its label satisfies
\texttt{label.strip().lower() == "palm"}; each palm shape, whatever its
\texttt{shape\_type}, is converted to its axis-aligned envelope
$(\min x, \min y, \max x, \max y)$. This yields 5{,}853 palm boxes
(3{,}306 rotation, 2{,}544 rectangle, 3 point shapes). Point shapes produce
zero-area boxes that can never be matched; they are retained.
% Source: src/preprocessing/gt_palm_bboxes.py (is_palm_label,
%         EVALUATION_PROTOCOL_VERSION="v2"); docs/EVALUATION_PROTOCOL.md §1
The annotations are not documented as exhaustive, so a real palm may lack a
LabelMe box.
\TODO{Who produced the LabelMe annotations, under what instructions, and whether exhaustiveness was intended.}

\paragraph{Protocol v1 $\rightarrow$ v2.} Until 2026-09-27 the label rule was
the case-sensitive \texttt{label == "palm"}. Seventy annotation files label palms
\texttt{"Palm"}; v1 therefore dropped 486 boxes. Protocol~v2 corrected only the
label rule; all stored predictions were re-scored without rerunning inference.
All results in this thesis use Protocol~v2.
% Source: docs/EVALUATION_PROTOCOL.md §7; docs/EXPERIMENT_RESULTS_CANONICAL.md §2

% ---------------------------------------------------------------------
\subsection{IoU-Based Annotation Matching}
\label{sec:method:iou}
% ---------------------------------------------------------------------
Matching is performed independently per image. Pairwise IoU is computed between
every detection and every LabelMe palm box; candidate pairs are sorted by
descending IoU and matched greedily one-to-one (each detection and each box is
used at most once). A pair is accepted iff $\mathrm{IoU} \geq 0.5$.
% Source: src/evaluation/gt_matching.py; docs/EVALUATION_PROTOCOL.md §2
A matched detection is labeled \emph{\matched{}} (legacy notation: GT+); an
unmatched detection is labeled \emph{\unmatched{}} (legacy: GT$-$). A detection
can be \unmatched{} because (a) no LabelMe box exists for the object, (b) its
best LabelMe box overlaps it with $\mathrm{IoU} < 0.5$, or (c) a
better-overlapping duplicate detection already took that box.
% Source: docs/SEMANTIC_VALIDITY_AUDIT.md §A; EVALUATION_PROTOCOL.md §2
The cohort contains 5{,}109 \matched{} and 638 \unmatched{} detections; an
always-\R{} verifier therefore attains alignment accuracy 0.889.
% Source: docs/EVALUATION_PROTOCOL.md §1.3

\paragraph{Terminology.} Throughout this thesis, ``\unmatched{}'' denotes an
IoU-matching outcome. It is \emph{not} synonymous with ``non-palm'' or
``detector false positive''; whether an \unmatched{} detection is semantically a
palm is a separate question answered only by the human audit
(\cref{sec:method:semantic}).

% ---------------------------------------------------------------------
\subsection{Input Conditions \Acond{1}--\Acond{5}}
\label{sec:method:conditions}
% ---------------------------------------------------------------------
Each sample is rendered under five conditions that vary the image and the
detector metadata included in the prompt (\cref{tab:conditions}). All images
derive from the same raw patch and box.
% Source: src/preprocessing/verification_overlay.py;
%         src/preprocessing/ablation_verification_images.py;
%         src/prompts/ablation_verification_prompts.py

\begin{itemize}
  \item \textbf{Overlay} (\Acond{1}--\Acond{3}, left panel of \Acond{4}): the
    full patch with every pixel outside the box multiplied by 0.55, the box
    region at original brightness, and a green rectangle (BGR $(0,200,0)$,
    2~px) drawn on the box.
  \item \textbf{\Acond{4} dual panel}: a canvas twice the patch width; left
    panel is the overlay, right panel is the unpadded box region cropped
    \emph{from the overlay} and resized to the patch size (aspect ratio not
    preserved).
  \item \textbf{\Acond{5} crop only}: the box expanded by 15~px on each side
    (clamped to the image), cropped from the \emph{raw} patch (no dimming, no
    rectangle), and resized to $512\times512$ (aspect ratio not preserved).
\end{itemize}

\begin{table}[t]
  \centering
  \small
  \caption{Input conditions as implemented. Metadata are inserted as text into the prompt.}
  \label{tab:conditions}
  \begin{tabularx}{\linewidth}{@{}l l X@{}}
    \toprule
    \textbf{Code} & \textbf{Image} & \textbf{Prompt metadata} \\
    \midrule
    \Acond{1} & Overlay & None (prompt instructs not to use confidence or geometry) \\
    \Acond{2} & Overlay & YOLO confidence \\
    \Acond{3} & Overlay & YOLO confidence; box width, height, area; aspect ratio (w/h) \\
    \Acond{4} & Overlay + enlarged crop (two panels) & YOLO confidence \\
    \Acond{5} & Enlarged padded crop only & YOLO confidence \\
    \bottomrule
  \end{tabularx}
\end{table}

\paragraph{Discrepancy with design documents.} \texttt{README.md} and
\texttt{docs/ABLATION\_STUDY.md} list ``center coordinates'' among the \Acond{3}
metadata. The implemented prompt builder does not include them and instead
instructs ``Do not use center coordinates or raw bbox coordinates''; the stored
\Acond{3} prompt files confirm this. This thesis reports the implemented
condition.

\paragraph{Confounds in the design.} (i) \Acond{1}, \Acond{2} and \Acond{3}
share an identical image, so their differences are purely textual. (ii)
\Acond{1}$\rightarrow$\Acond{4} and \Acond{1}$\rightarrow$\Acond{5} change both
the image and the text (confidence is added); the image-only contrasts are
\Acond{2}$\rightarrow$\Acond{4} and \Acond{2}$\rightarrow$\Acond{5}. (iii)
\Acond{4} and \Acond{5} differ in more than surrounding context: the \Acond{4}
crop comes from the overlay (dimmed outside, rectangle visible at the box
border) without padding, whereas the \Acond{5} crop comes from the raw image with
15~px padding and a fixed $512\times512$ output. (iv) The condition name (e.g.,
\texttt{A3\_overlay\_confidence\_geometry}) and the sample ID appear in every
prompt (\cref{sec:method:prompt}).

\begin{figure}[t]
  \centering
  % TODO: AUTO-GENERATE FROM RESULTS
  % Candidate source: docs/assets/readme/ablation_inputs_light.png — BUT that
  % illustration uses a LabelMe box, not a YOLO candidate (README caption).
  % Regenerate with a real YOLO candidate from outputs/verification_ablation_5747/.
  \FIGPLACEHOLDER{One YOLO candidate rendered under A1/A2/A3 (identical image), A4 and A5, taken from \texttt{outputs/verification\_ablation\_5747/}.}
  \caption{The five input conditions for one candidate detection.}
  \label{fig:conditions}
\end{figure}

% ---------------------------------------------------------------------
\subsection{Prompt Protocol}
\label{sec:method:prompt}
% ---------------------------------------------------------------------
Prompts are generated once per sample and condition by
\texttt{build\_ablation\_verification\_prompt()} and stored as text files;
adapters read these files and never rebuild prompts. All models therefore
receive byte-identical prompt text, wrapped in their own chat template.
% Source: src/prompts/ablation_verification_prompts.py;
%         scripts/pipeline/build_ablation_verification_prompts.py;
%         docs/FRAMEWORK_FREEZE.md "Prompt builder"
Every prompt contains, in order: a role statement (expert reviewer; the task is
to verify, \emph{not} to detect new palms); general instructions (focus on the
highlighted box; use surrounding context only if useful; YOLO confidence is
auxiliary and must not be the primary basis); a list of wild-palm visual
characteristics; a condition-specific image description; the decision
definitions; reasoning requirements; the condition-specific metadata block and
instruction (\cref{tab:conditions}); the condition name and sample ID; and a
request to return only JSON in a fixed template. The full \Acond{1} prompt is
reproduced in \cref{app:prompt}.

Two properties of the shared text are relevant for interpretation. First, the
general instructions mention the YOLO confidence score in every condition,
including \Acond{1}, where no score is given. Second, the \Acond{1} instruction
``Judge only from the highlighted bounding box in the image'' differs from the
general instruction to use surrounding context when useful.
\TODO{Decide whether to discuss these as threats to validity in \cref{sec:discussion}.}

% ---------------------------------------------------------------------
\subsection{Output Schema and Parsing}
\label{sec:method:schema}
% ---------------------------------------------------------------------
The model is asked to return
\begin{verbatim}
{
  "decision": "Reliable | Uncertain | Unreliable",
  "confidence_reasoning": "",
  "visual_reasoning": ""
}
\end{verbatim}
with the following decision definitions (verbatim from the prompt):
\begin{description}[leftmargin=1.5em]
  \item[\R{}] ``The highlighted bounding box clearly contains a valid wild palm.''
  \item[\U{}] ``Evidence is ambiguous due to occlusion, overlap, low image quality, partial visibility, or confusing vegetation.''
  \item[\Ur{}] ``The highlighted object is not a palm or the detection is clearly incorrect.''
\end{description}
Note that \Ur{} covers both semantic rejection (not a palm) and localization
rejection (``clearly incorrect''); it is not a pure non-palm verdict.

Parsing is shared: model-specific raw-text cleanup
(\texttt{normalize\_raw\_response}; identity for all models except the
Ministral~3 keys), then JSON extraction (direct parse, then a fenced
\texttt{```json} block, then the first \texttt{\{...\}} span), then decision
normalization, which accepts exactly the three labels case-insensitively. Any
other output is recorded as a \texttt{parse\_error} with an empty decision and is
never coerced into a label.
% Source: src/lvm/parsers/base.py; src/lvm/response_schema.py::parse_json_response;
%         src/lvm/parsers/cleanup.py; FRAMEWORK_FREEZE.md amendment 2026-10-01
Each prediction is stored as one JSON record per sample containing at least the
decision, both reasoning fields, the raw response, the parsed response, runtime,
parse/inference error fields, and a UTC timestamp.

% ---------------------------------------------------------------------
\subsection{Models and Execution Protocol}
\label{sec:method:models}
% ---------------------------------------------------------------------
Each model is integrated through a verifier and an adapter that use the model's
native processor and chat template, a per-model YAML configuration, and a
registry entry. The shared runner handles job iteration, resume (keyed by
\texttt{sample\_id} within a results directory) and output persistence; each
job loads the model once and runs one condition. All configurations use greedy
decoding (\texttt{do\_sample=False}) with \texttt{max\_new\_tokens=512};
Gemma~4 runs with \texttt{enable\_thinking=False}.
% Source: docs/FRAMEWORK_FREEZE.md; configs/models/*.yaml;
%         docs/EXPERIMENT_STATUS_CANONICAL.md §3
\TODO{Per-model image preprocessing (resolution limits, tiling, max pixels) applied by each native processor; these differ across families and are not summarized in the repository.}

\Cref{tab:models} lists every model integrated in the repository and its
evaluation status as of this draft.

\begin{table}[t]
  \centering
  \scriptsize
  \caption{Integrated models and evaluation status (N = 5{,}747 unless stated). Status as recorded in \texttt{docs/EXPERIMENT\_STATUS\_CANONICAL.md}, \texttt{docs/MODEL\_EXPANSION\_CAMPAIGN.md} and the Slurm queue on 2026-10-02. \TODO{Refresh before submission.}}
  \label{tab:models}
  \begin{tabularx}{\linewidth}{@{}l l l X@{}}
    \toprule
    \textbf{Family} & \textbf{Checkpoint / size} & \textbf{Registry key} & \textbf{Status} \\
    \midrule
    Qwen2.5-VL & 7B-Instruct & \texttt{qwen2\_5\_vl} & A1--A5 complete (legacy launch path; A5 resumed) \\
               & 3B          & \texttt{qwen2\_5\_vl\_3b} & A1--A5 queued \\
               & 32B         & \texttt{qwen2\_5\_vl\_32b} & qualification probe pending (H200) \\
    Qwen3-VL   & 8B-Instruct & \texttt{qwen3\_vl} & A1--A5 complete \\
               & 4B          & \texttt{qwen3\_vl\_4b} & A1--A5 predictions and v2 evaluation present; not yet in canonical docs \\
               & 2B          & \texttt{qwen3\_vl\_2b} & A1, A2 running; A3--A5 queued \\
               & 32B         & \texttt{qwen3\_vl\_32b} & qualification probe pending (H200) \\
    InternVL3.5-HF & 8B      & \texttt{internvl3\_5\_hf} & A1--A5 complete \\
               & 2B, 4B, 14B & \texttt{internvl3\_5\_hf\_\{2b,4b,14b\}} & A1--A5 queued \\
    GLM        & GLM-4.6V-Flash \TODO{size} & \texttt{glm\_4\_6v\_flash} & A1--A5 complete \\
    Phi-4      & multimodal-instruct \TODO{size} & \texttt{phi4\_multimodal} & A1--A5 complete \\
    Gemma 4    & 12B IT      & \texttt{gemma4} & A1--A5 complete (skipped qualification ladder) \\
    MiniCPM-V  & 4.5         & \texttt{minicpm\_v4\_5} & A1 complete (H200); A2--A5 queued (H200) \\
    Molmo2     & 8B          & \texttt{molmo2\_8b} & A1 complete (H200); A2--A5 queued (H200) \\
    LLaVA-OneVision & Qwen2-7B OV & \texttt{llava} & A1 @1{,}000 only \\
    Gemma 3    & 12B IT      & \texttt{gemma} & A1 @1{,}000 only \\
    InternVL3  & 8B-Instruct & \texttt{internvl3} & stopped at qualification (parse failures) \\
    Llama 3.2 Vision & 11B-Instruct & \texttt{llama3\_2\_11b\_vision} & engineering check only; held \\
    Ministral 3 & 3B, 8B, 14B & \texttt{ministral3\_\{3b,8b,14b\}} & qualification probes pending \\
    \bottomrule
  \end{tabularx}
\end{table}

\paragraph{Hardware and environments.} Full-scale runs used single NVIDIA L40S
GPUs, except MiniCPM-V-4.5 and Molmo2-8B (H200; their A2--A5 runs are also
scheduled on H200 for parity), and the 32B probes (H200). Several families run
in isolated Python environments with different \texttt{transformers}
versions. Attention backend is SDPA for the scaling runs; Phi-4 uses eager
attention.
% Source: EXPERIMENT_STATUS_CANONICAL.md §3; MODEL_EXPANSION_CAMPAIGN.md;
%         configs/models/phi4_multimodal.yaml
\TODO{Complete per-model environment table (Python, \texttt{transformers}, \texttt{torch}, dtype, attention, revision) in \cref{app:repro}. Revision of the Qwen2.5-VL-7B anchor was not recorded.}

\paragraph{Qualification.} Models passed through qualification runs before
full-scale inference, but the ladder was \emph{not} applied uniformly
(\texttt{EXPERIMENT\_STATUS\_CANONICAL.md} §5). The original gates
(balanced-100) used alignment-based thresholds; the scaling campaign uses a
GT-independent paired probe (the first 100 samples under all five conditions)
with only technical hard gates (record completeness, zero inference errors,
$\geq 95\%$ parseable outputs per condition, valid labels, pinned revision, one
attention backend).
% Source: docs/MODEL_EXPANSION_CAMPAIGN.md "Technical qualification"
\TODO{Decide how to report gate deviations (Gemma~4 skipped Stage~0/balanced-100/A1@1000; InternVL3.5 skipped A1@1000).}

\paragraph{Run integrity.} For every completed full-scale run used here, the
repository records exactly 5{,}747 prediction files per condition, all with
status \texttt{ok}, zero parse errors and zero inference errors, and decision
counts that match the evaluation files.
% Source: EXPERIMENT_STATUS_CANONICAL.md §2; A1_A5_BEHAVIOR_ANALYSIS.md §2
\VERIFY{Re-run the integrity audit for every scaling-campaign run (Qwen3-VL-4B and later) before including it.}

% ---------------------------------------------------------------------
\subsection{Evaluation}
\label{sec:method:eval}
% ---------------------------------------------------------------------
Results are reported at three layers that measure different constructs.

\paragraph{Layer 1 — Behavior (reference-free; primary).} Computed from stored
decisions only: the distribution of \R{}/\U{}/\Ur{} over all $N$ samples;
coverage $=(\R+\Ur)/N$; abstention $=\U/N$; per-detection $3\times3$ transition
matrices between conditions on identical sample IDs; the share of decisions that
change; direction of change along the order
\R{}$\rightarrow$\U{}$\rightarrow$\Ur{} (``toward rejection'' versus ``toward
acceptance''; neither direction is better); and per-detection stability (number
of distinct labels over \Acond{1}--\Acond{5}).
% Source: docs/A1_A5_BEHAVIOR_ANALYSIS.md §3

\paragraph{Layer 2 — Annotation alignment (Protocol v2; secondary).} \R{} is a
positive prediction, \Ur{} a negative prediction, and \U{} is excluded. Against
the matching label:
\begin{center}
\small
\begin{tabular}{lcc}
  \toprule
  & \matched{} & \unmatched{} \\
  \midrule
  \R{}  & TP & FP \\
  \Ur{} & FN & TN \\
  \U{}  & excluded & excluded \\
  \bottomrule
\end{tabular}
\end{center}
Accuracy, precision, sensitivity (recall), \alignspec{} $=\mathrm{TN}/(\mathrm{TN}+\mathrm{FP})$,
F1 and balanced accuracy $=(\text{sensitivity}+\text{specificity})/2$ are
computed on decided samples only. Because the class prior is high (0.889) and
\U{} is excluded, accuracy and F1 are always reported with \alignspec{},
balanced accuracy and coverage. In this thesis TP/FP/TN/FN and
specificity are \emph{alignment} quantities: ``FP'' means \R{} on a
\unmatched{} detection, not acceptance of a non-palm.
% Source: docs/EVALUATION_PROTOCOL.md §3--4; scripts/compute_verification_metrics.py

\paragraph{Layer 3 — Semantic audit.} See \cref{sec:method:semantic}.

\TODO{Specify uncertainty quantification for Layers 1--2 (e.g., image-clustered bootstrap CIs; see \texttt{ANALYSIS\_PLAN.md}). The existing detection-property analysis uses a source-image bootstrap; reuse its design if appropriate.}

% ---------------------------------------------------------------------
\subsection{Semantic Review Procedure}
\label{sec:method:semantic}
% ---------------------------------------------------------------------
\paragraph{Unmatched-detection audit.} All 638 \unmatched{} detections of the
cohort were visually reviewed by one human reviewer, who saw only the raw patch
and the target box. No LabelMe annotation, YOLO confidence, IoU, other detection
or VLM output was shown. Labels were \texttt{palm}, \texttt{non\_palm} or
\texttt{ambiguous}. Labeling took place on 2026-10-01 through a dedicated blind
review application.
% Source: docs/SEMANTIC_VALIDITY_AUDIT.md §B; scripts/semantic_review_app.py
\TODO{Reviewer identity/expertise and written labeling criteria (``criteria were standardized'' after a familiarization round, per the audit record; locate the criteria).}

\paragraph{Lower-confidence pilot.} To test whether non-palm candidates exist
\emph{below} the cohort threshold, 400 YOLO detections with confidence in
$[0.10, 0.50)$ were sampled, 100 per 0.1-wide bin (with patch and frame caps;
selection seed 20261001; review-order seed 20261002), and reviewed blind in
shuffled order. These detections are outside the verification cohort and have
no VLM predictions. Because bins were sampled equally, raw pooled shares are not
prevalence estimates.
% Source: docs/SEMANTIC_VALIDITY_AUDIT.md §C--D

\paragraph{Limitations of the audit (recorded in the repository).} Single
reviewer (no inter-rater agreement); localization quality not assessed (a
\texttt{palm} label does not mean the box is well localized, whereas \Ur{} also
covers ``clearly incorrect'' detections); \matched{} detections were not
reviewed; no matched controls were interleaved; the review was fast (638 items in
about 13 minutes).
% Source: docs/SEMANTIC_VALIDITY_AUDIT.md §F

\paragraph{Superseded pipeline.} An earlier semantic-GT pipeline
(\texttt{src/evaluation/semantic\_gt.py},
\texttt{outputs/semantic\_gt\_evaluation/}) predates the official review and is
marked superseded; it is not used for any result in this thesis.
\VERIFY{The IoU-sensitivity outputs are written under \texttt{outputs/semantic\_gt\_evaluation/iou\_sensitivity/}; confirm they do not depend on superseded semantic labels before citing them.}
